% Cote 119: offprints, reports, courses, CV and letters (typed and printed part): archive cote 119, site page 246 (PDF page 247).
% typescript. Draft compared: none.
% AI-transcribed by Claude Opus 5.5 (claude-opus-5-5, adaptive thinking, effort high), Message Batches API, batch r119_full (msgbatch_01DwZNm5Y7YsKBWyAN4sFs8w), 2026-11-13; not yet checked by eye.
\scanpage{119/246}
\mspage{6}

le plus souvent, un bon enseignant, et cela d'autant plus qu'il expose des
parties de la science qui sont d'un niveau théorique plus élevé et plus
proche des confins du connu. Il pourra être également bon enseignant au
niveau le plus élémentaire, traitant des bases même de la science, \uline{à condition}
qu'il s'adresse à un auditoire non pas de futurs utilisateurs sans plus,
mais à des étudiants qui se destinent à un approfondissement poussé de cette
science. Il n'est pas rare, sous cette réserve, que des savants de premier
ordre soient en même temps des enseignants émérites, et qu'ils aient pour
l'enseignement une passion qui augmente avec les années. Il arrive également,
mais c'est l'exception, qu'un chercheur de grande valeur soit un excécrable
enseignant : l'enseignement lui est une corvée qu'il accomplit à contre-coeur,
et il est autant à plaindre que ses étudiants, qui dès le début de l'année
académique perdent l'espoir de jamais pouvoir s\over{\unsure{'y}}{\unsure{'y}} retrouver dans les cours
qu'il dispense, touffus et laborieux quand ils ne sont obscurs et sibyllins.
Ils garderont sans doute le plus morne souvenir d'un homme dont ils n'ont
aucune raison de suspecter que c'est une des lumières de la science qu'il
leur enseigne si mal ! Personne, et lui-même moins que tout autre, ne con-%
testera que cet homme n'aurait jamais dû enseigner.

Le test le plus caractéristique pour distinguer le chercheur mathé-%
maticien de toute autre espèce d'individu, c'est que lorsqu'il rencontre
un congénère, que ce soit autour d'une table, dans un métro, dans un congrès
ou n'importe où que ce soit, il se met à parler aussitôt non de politique,
ni du temps, ni de ses élèves, ni de ses collègues, ni de ses patrons (quand
il en a), mais de problèmes mathématiques. L'enseignement l'intéresse ou
même le passionne par les problèmes techniques d'exposition qu'il soulève,
variables à l'infini suivant le sujet enseigné et les connaissances supposées
acquises, qui mettent en jeu de mille façons imprévues son ingéniosité et
sa virtuosité,et non en fonction de la personnalité propre de ses divers
auditeurs. On peut rapporter à ce sujet une anecdote bien connue (dont il
n'est pas question de garantir l'authenticité !), selon laquelle le mathé-%
maticien C.L. Siegel, connu également pour son amour de l'enseignement,
un jour que l'unique auditeur qui lui restait (il faut croire que c'était un

\margin{246}
